\section{Conclusion}
\label{sec:conclusion}

We began by asking why a symmetry-enforcing architecture fails at the very scale where its inductive bias should matter most. Our answer: it does not fail because the structural inductive bias is flawed---it fails because the bias requires data to activate. Below approximately 150--250 training models on the CIFAR-10 CNN zoo, the GNN-NFN graph encoder lacks sufficient topological diversity to benefit from permutation equivariance. Above this threshold, the structural constraint reduces the effective hypothesis space and produces a large sample efficiency advantage over a plain flat-MLP.

Our contributions are: (1) a large efficiency advantage---the first controlled measurement on shared ModelZooDataset splits, showing GNN-NFN reaches 90\% of its peak performance while flat-MLP requires the full training dataset to reach 90\% of its own peak R$^2{=}0.886$; (2) a novel data-regime crossover (PermAug $>$ GNN-NFN at $N{=}100$; GNN-NFN $>$ PermAug at $N{\geq}250$) revealing a minimum-data threshold of ${\sim}150$--250 models; (3) floating-point-precision equivariance verification on trained models grounding the efficiency claim mechanistically; and (4) a reusable evaluation protocol for future encoder comparisons.

Three directions follow directly from our results: multi-seed replication of the $N{=}100$ crossover (low cost, high priority); GNN-NFN learning-rate ablation at $N{=}100$ to distinguish underfitting from a fundamental equivariance limit; and extension to the MNIST MLP zoo with DWSNets to assess generality across zoo types.

As model zoos grow and equivariant encoder designs mature, the minimum-data threshold we identify may shrink---but until it does, practitioners with small zoos should start with augmented simplicity, not equivariant complexity. The question of when structural inductive bias pays off is not merely architectural; it is fundamentally a question about data.
